\chapter{Type 2 Diabetes Mellitus}
\index{Type 2 Diabetes Mellitus}
\label{sec:Type 2 Diabetes Mellitus}

\section{Overview}
Type 2 diabetes mellitus (T2DM) is the most common form of diabetes, accounting for 90--95\% of cases. It is marked by insulin resistance in peripheral tissues (muscle, liver, adipose tissue) combined with progressive beta-cell dysfunction and relative insulin deficiency. This metabolic disorder involves dysregulation of normal bodily processes, often requiring ongoing monitoring and management. It is increasingly common due to lifestyle and dietary factors. Metabolic disorders involve disruption of normal biochemical processes essential for health. These conditions may result from enzyme deficiencies, hormonal imbalances, or lifestyle factors. Many metabolic disorders are chronic and require ongoing management. Lifestyle modifications are often the cornerstone of treatment. This condition affects a significant number of people worldwide and can range from mild to severe in presentation. Early recognition and appropriate management are essential for optimal outcomes. A thorough understanding of the underlying mechanisms guides treatment decisions. Patient education and self-management play important roles in long-term care.
\section{Causes}
The Disease mechanism involves the "ominous octet"---eight mechanisms contributing to hyperglycemia: decreased insulin secretion, increased glucagon secretion, increased liver glucose production, decreased muscle glucose uptake, increased lipolysis, decreased incretin effect, increased kidney glucose reabsorption, and neurotransmitter dysfunction. Metabolic disorders arise from dysregulation of normal biochemical pathways. This may result from enzyme deficiencies, hormonal imbalances, or lifestyle factors that overwhelm normal homeostatic mechanisms. The underlying mechanisms involve complex interactions between genetic predisposition and environmental factors. Disruption of normal physiological processes leads to the characteristic features of the condition. Multiple contributing factors may act synergistically to trigger or worsen the condition.
\section{Risk Factors}
\begin{itemize}[noitemsep]
  \item Risk factors include obesity (particularly visceral adiposity), physical inactivity, family history, age >45, ethnicity, history of gestational diabetes, and polycystic ovary syndrome. Your outlook depends on glycemic control and Treatment for cardiovascular risk factors
  \item Cardiovascular benefit is demonstrated with SGLT2 inhibitors and GLP-1 receptor agonists.
  \item Family history
  \item Obesity and sedentary lifestyle
  \item Poor dietary habits
  \item Advancing age
  \item Hormonal imbalances
  \item Certain medications
  \item Genetic predisposition and family history
  \item Lifestyle factors (diet, exercise, smoking, alcohol)
  \item Environmental exposures
  \item Underlying medical conditions
  \item Medication use
\end{itemize}
\section{Signs and Symptoms}
\begin{itemize}[noitemsep]
  \item Pain or discomfort in the affected area
  \item Functional impairment affecting daily activities
  \item Changes in appearance or function of affected tissues
  \item Systemic symptoms may include fatigue, fever, or weight changes
  \item Symptoms may develop gradually or appear suddenly
\end{itemize}
\section{Progression and Stages}
Symptoms of advanced disease include polyuria, polydipsia, and weight loss. The condition may progress through different stages of severity over time. Early stages often involve mild or subtle symptoms that may be overlooked. Moderate stages involve more noticeable symptoms affecting daily function. Advanced stages may involve significant impairment and complications.
\section{Complications}
Complications include microvascular disease (retinopathy, nephropathy, neuropathy) and macrovascular disease (coronary artery disease, stroke, peripheral arterial disease).

\begin{itemize}[noitemsep]
  \item Disease progression leading to worsening symptoms
  \item Secondary infections or injuries
  \item Side effects of treatment
  \item Reduced quality of life
  \item Emotional and psychological impact
\end{itemize}
\section{Diagnosis}
Doctors diagnose this based on fasting plasma glucose, HbA1c, or OGTT.

\begin{itemize}[noitemsep]
  \item Comprehensive medical history and physical examination
  \item Laboratory tests to assess organ function
  \item Imaging studies to visualize affected structures
  \item Specialized testing depending on the affected system
  \item Consultation with relevant specialists
\end{itemize}
\section{Prevention}
Many patients are asymptomatic and diagnosed on routine screening. Management follows a stepwise approach: lifestyle modification (diet, exercise, weight loss of 5--10\%), metformin as first-line pharmacotherapy, followed by SGLT2 inhibitors, GLP-1 receptor agonists, DPP-4 inhibitors, thiazolidinediones, sulfonylureas, and insulin as needed.

\begin{itemize}[noitemsep]
  \item Maintain a healthy lifestyle with balanced diet and exercise
  \item Avoid known risk factors
  \item Attend regular medical check-ups for early detection
  \item Follow recommended screening guidelines
\end{itemize}
\section{Treatment Options}
Dietary modifications and nutritional counseling Pharmacotherapy to correct metabolic abnormalities Hormone replacement therapy when indicated Regular monitoring of metabolic parameters Weight management programs Exercise prescription Management of associated conditions Patient education for self-management

\begin{itemize}[noitemsep]
  \item Medications to manage symptoms and address underlying mechanisms
  \item Lifestyle modifications and self-care measures
  \item Physical therapy and rehabilitation
  \item Surgical intervention when indicated
  \item Regular monitoring and follow-up care
\end{itemize}
\section{Living with It}
\begin{itemize}[noitemsep]
  \item Follow your treatment plan as prescribed
  \item Maintain regular follow-up with your healthcare provider
  \item Make necessary lifestyle adjustments
  \item Seek support from family, friends, or support groups
  \item Stay informed about your condition
\end{itemize}
